% !TeX root = ../papers/nc-p3.tex
\section{指定した判別直線と二次曲線を持つ階数十の正則変形}
\label{sec:marked-prescribed-discriminant-as}
\begin{theorem}\label{thm:marked-prescribed-discriminant-as}
三次元内積空間 $V$ において任意の非零跡零対称行列 $N$ と任意の射影直線 $\ell(v)=0$ を指定する。
単純印付き $q=3$, $g=4$ の四窓に、判別三次式が定数倍を除き
\[
                       \ell(v)\,v^tNv
\]
となる十九次元の既約族があり、その幾何学的一般点は四辺有限、階数10の標準Koszul左右AS正則窓である。
二次曲線の階数、$N$ のJordan型、直線との接触や成分の一致に例外はない。
$[N],[\ell]$ の射影直交群内の同時安定化群の次元を $d$ とすると、
この族の添字保持同型類の像の次元は $9-d$ である。
また $N,\ell$ をともに一般に動かした場合には、同型類の像の次元は12になる。
\end{theorem}
\begin{proof}[混合正則点でのエタールな復元]
$S=\mathrm{Sym}_0(V)$, $D=N^\perp\subset S$ とする。順序付き三列 $G$ を $D$ 内で選び、
三平面 $G'$ の写像 $\bigwedge^2G'\to\mathfrak{so}(V)$ が可逆となる開集合を取る。
実際にここで使える正則な基点は後で全 $N$ について示す。
$H\subset G'$ を、$[H_1,H_2]$ の跡対合による一次式が $\ell$ を定める一意な二平面とする。
これは三次元の二ベクトルが全て可分解であるためで、$G$ を動かすと代数的に動く。
$H\subset F'\subset D$ となる三平面を選ぶ自由度は $\mathbf P^1$ である。
その順序付き基底 $F_i$ に課す式は $\sum_i[G_i,F_i]=0$ である。
これがちょうど指定三交換子の跡零条件である。
$F'=G'$ では写像 $(F_i)\mapsto\sum_i[G_i,F_i]$ が全射なので、
近傍の三式の核は階数6の束となる。その中で基底となることは開条件で
$F=-G$ が使える。得られた滑らかな既約基底の次元は $12+1+6=19$ であり、
$F'\ne G'$ なら $G'\cap F'=H$, $G'+F'=D$ である。

$F=-G$ では三つの指定値は全て交代行列で、その三列は独立である。
ここでいう指定値は $b,d$ を固定し $a,c$ の行列としたものである。
外積行列の二交換子は二列の外積なので、$Q_v=[Lv]_\times$ と取れる。
実際 $\operatorname{cof}(L)=C$ は可逆三行列 $C$ に対し
$L=\sqrt{\det C}\,C^{-t}$ で解ける。$C$ は指定値の外積ベクトルを並べた行列である。
三交換子写像 $\mathfrak{sl}_3^3\to\mathfrak{sl}_3^3$ はこの点でエタールである。
基底を取り直して $Q_i=-[e_i]_\times$ で確かめれば十分で、
微分 $(T_i)\mapsto([T_i,Q_j]+[Q_i,T_j])_{ij=23,31,12}$ の行列式は
本文末の再現計算の基底順で $6912$ である。
この写像を上の十九次元基底に沿って引き戻すと、指定解を通る唯一の局所成分は滑らかで基底に支配的である。
その既約開部分が所要のパラメータ族を与える。一般点での三次式は
命題\ref{prop:marked-mixed-discriminant}による。

次の段で混合点の辺の有限性を確認すれば、定理\ref{thm:mixed-point-frobenius-algebras}により
そこで全次数正則性と階数10が成立する。符号 $F=-G$ は偶二補空間の基底を
互いに逆の平方根で調整して $F=G$ にでき、巡回順を替えれば同定理の交代・対称表示である。
二つの隣接切断の階数6は交換子式による。一文字階数、三・四次交叉、五次双対消滅は
混合点の近傍で標準値を保つ。下界は与えた四テンソルとその縮約、
上界は関係挿入行列の階数による。点辺の有限性も射影族で開であり、
星印は全族で六関係の一次部が独立な単純点を与える。
反対階数は基点で10、近傍で片側対称性により10以下である。
この近傍の有限Frobenius双対の相対bar複体を内部次数ごとに取ると、
非対角ホモロジーの消滅は各次数で基点を含む開条件である。
幾何学的一般点では全次数で消滅し、双対の完全対合により左右AS条件も従う
（定理\ref{thm:reversible-component}の一般点の証明）。
これは十九次元族の全閉点での一様な正則性判定を与える主張ではない。
\end{proof}

\begin{proof}[各跡法線での有限混合点]
半単純な $N$ は直交対角化できる。$G_i=E_{jk}+E_{kj}$（巡回順）を選べば $G_i\in D$ で、
$F=-G,Q_i=-[e_i]_\times$ は定理\ref{thm:twenty-point-distinct-companion}の二十点例を
一位置巡回した星印での表示である。実際同定理の $B_+^t B_-$ の四因子を $(v,x,y,u)$ の順で読むと、
星印一個の係数は順に $D_i,-D_i,-[e_i]_\times,-[e_i]_\times$ となる
（ここで $D_i=E_{jk}+E_{kj}$）。よってこの場合は一般点でも点スキームは被約二十点となる。

非半単純なら、相似と非零スカラー倍でJordan型は平方零階数1、正則冪零、非零重複固有値の二次ブロックの
三通りである。自己随伴行列の相似は計量を保つ相似にできる。標準計量で二つの対称行列を
$T$ で共役すれば $T^tT$ が元の行列と交換するので、その可逆な多項式平方根で $T$ を補正すればよい。
従って反対角単位行列 $J$ を計量にし
\[
 H(a,b,c,d,e)=\begin{pmatrix}a&b&c\\d&-2a&b\\e&d&a\end{pmatrix}
\]
とした次の三代表で尽きる。表では $N$ と $G_i$ を引数の五組で記す。
\[
\begin{array}{c|c|c|c}
N&G_1&G_2&G_3\\ \hline
(0,0,1,0,0)&(-2,-1,2,-2,0)&(-2,-2,1,-2,0)&(2,0,2,-2,0)\\
(0,1,0,0,0)&(2,2,0,0,1)&(-2,-1,-2,0,-1)&(2,-2,1,0,0)\\
(1,0,1,0,0)&(2,12,6,-6,-12)&(0,0,12,-12,0)&(-1,12,0,-6,6)
\end{array}
\]
各々で $\operatorname{tr}(NG_i)=0$ で、$C$ の行列式は順に $-576,18,-1306368$。
ここで $C$ は巡回 $bd$ に対する $(J[G_a,G_c])_{bd}$ を $ac$ 交代行列として外積ベクトルにした三列である。
$L_0=\operatorname{adj}(C)^t$ は $L$ の非零倍であり、星座標の倍率でこの差を吸収できる。
辺の係数行列を点 $(z,v)$ で評価すると、列変数を $(c_*,c)$ として
\[
          \begin{pmatrix}0&zI+[L_0v]_\times\\v&-G_1v\quad -G_2v\quad -G_3v\end{pmatrix}
\]
となる。階数は全点で3以上である。$v\ne0$ なら上の三行の階数は2以上
（非零外積行列は巡回）、下に第一列の非零成分がある。$v=0$ なら上の三行が独立である。
平面 $z=v_1+2v_2+3v_3$ 上、十五個の四次小行列式は三変数四次式の全空間を張る。
実際 $(v_1,v_2,v_3)=(i,j,1)$, $i,j\ge0$, $i+j\le4$ の十五点での値の行列式は
$101$ を法として順に $53,38,82$ である。従って標数零で独立であり、
この平面上に退化点はない。射影的な退化軌跡は有限、各ファイバーは一点以下である。
混合点の逆辺も対称行列の転置と交代行列の符号反転だけが違うので、
星座標の符号変更および二位置の交換で同じ有限性を得る。
\path{computations/verify_prescribed_discriminant_seeds.py} が小行列、微分と次の安定化群の計算を再現する。
\end{proof}

\begin{proof}[同型類の次元]
定理\ref{thm:marked-point-commutators}の残余群では、$Q=P$ の三独立な基底を
保つ条件が $BB^t=sI$ であり、自由な変換群は
$\mathrm{GL}(U)\times\mathrm{GO}(V)$ で次元13である。
$g=4$ なら法線と共通二平面の交換子方向はテンソルと印から固有なので、
両者の射影値を固定した場合の群の次元は $9+1+d$ である。
基点の $G'$ を保つ射影直交群は有限である。半単純な場合は対角平面の安定化群、
残る三例では $J$ 交代な $T$ に対する $[T,G_i]\bmod G'$ の係数行列が階数3である。
従って連結安定化群では $B$ がスカラーで、$G$ の三独立な順序付き基底を保つため $A=I$、
$F$ を保つ条件から $s=1$ となる。近傍でも安定化群は有限である。
連結な変換群は既約成分を保ち、開部分での軌道の次元は上の群次元に等しい。
印を忘れる選択肢も有限であるので $19-(10+d)=9-d$ を得る。
一般に動かす場合は $[N]\in\mathbf P^4$ と $[\ell]\in\mathbf P^2$ の六変数を基底に加えた
同じ束とエタールチャートを用いる。次元 $25-13=12$ が従う。
\end{proof}
